\documentclass[11pt]{article}

\usepackage[letterpaper,margin=1in]{geometry}
\usepackage[T1]{fontenc}
\usepackage[utf8]{inputenc}
\usepackage{lmodern}
\usepackage{microtype}
\usepackage{amsmath,amssymb,mathtools,bm}
\usepackage{xcolor,slashed,needspace}
\usepackage{enumitem}
\usepackage{hyperref}
\usepackage[most]{tcolorbox}
\usepackage{fancyhdr}
\usepackage{parskip}

\definecolor{qftblue}{HTML}{245A8D}
\definecolor{qftgray}{HTML}{F4F6F8}
\hypersetup{
  colorlinks=true,
  linkcolor=qftblue,
  urlcolor=qftblue,
  pdftitle={PHYS 662 Homework Set 5},
  pdfauthor={Nima Lashkari}
}

\pagestyle{fancy}
\fancyhf{}
\lhead{PHYS 662: Quantum Field Theory I}
\rhead{Homework Set 5}
\cfoot{\thepage}
\setlength{\headheight}{14pt}

\setlist[enumerate]{leftmargin=2.2em,itemsep=0.55em,topsep=0.5em}
\newlist{parts}{enumerate}{1}
\setlist[parts]{label=\textbf{(\alph*)},leftmargin=2.3em,itemsep=0.65em}

\newcommand{\dd}{\mathrm{d}}
\newcommand{\ii}{\mathrm{i}}
\newcommand{\R}{\mathbb{R}}
\newcommand{\C}{\mathbb{C}}
\newcommand{\Z}{\mathbb{Z}}
\newcommand{\Lag}{\mathcal L}
\newcommand{\tr}{\operatorname{tr}}
\newcommand{\Ad}{\operatorname{Ad}}

\newtcolorbox{conventionbox}[1][]{
  enhanced,breakable,
  colback=qftgray,colframe=qftblue,
  boxrule=0.8pt,arc=1.5mm,
  title={#1},fonttitle=\bfseries,coltitle=white
}
\newcommand{\problem}[2]{%
  \clearpage
  \section*{\color{qftblue}Problem #1: #2}
  \addcontentsline{toc}{section}{Problem #1: #2}
}
\newcommand{\lectureconnection}[1]{%
  {\small\textit{Lecture connection: #1.}\par}\medskip
}


\newcommand{\diag}{\operatorname{diag}}
\newcommand{\SO}{\mathrm{SO}}
\newcommand{\SU}{\mathrm{SU}}
\newcommand{\SL}{\mathrm{SL}}
\newcommand{\ISO}{\mathrm{ISO}}
\newcommand{\ddPi}{\dd\Pi}
\newcommand{\vac}{\lvert0\rangle}
\newcommand{\vbra}{\langle0\rvert}
\newcommand{\Dc}{\mathcal C}
\newcommand{\one}{\mathbf1}
\begin{document}
\begin{center}
{\scshape\Large PHYS 662 --- Quantum Field Theory I\par}
\vspace{0.25em}
{\Huge\bfseries\color{qftblue} Homework Set 5\par}
\vspace{0.25em}
{\large Lectures 15--17 \quad Fall 2026\par}
\end{center}
\vspace{0.5em}
\textbf{Name:} \rule{0.44\textwidth}{0.4pt}\hfill
\textbf{Due:} \rule{0.24\textwidth}{0.4pt}

\begin{conventionbox}[Instructions and conventions]
Show derivations and state distributional prescriptions. Set
\(\hbar=c=1\) and use the mostly-minus metric. Our notation is
\[
\eta=\diag(1,-1,-1,-1),\qquad p\cdot x=E_{\bm p}t-\bm p\cdot\bm x,
\qquad \dd\Pi_p=\frac{\dd^dp}{(2\pi)^d2E_{\bm p}},
\]
where \(E_{\bm p}=\sqrt{\bm p^2+m^2}\).
Take \(m>0\) unless a massless limit is requested.

For the real scalar field, use
\[
 D^+(z)=\langle0|\phi(x)\phi(y)|0\rangle,\qquad
 \mathcal K(z)=D^+(z)-D^+(-z)=[\phi(x),\phi(y)],\quad z=x-y.
\]
Set \(L=\Box+m^2\), \(G_R=i\theta(t)\mathcal K\),
\(G_A=-i\theta(-t)\mathcal K\), and
\(D_F=\langle0|T\phi(x)\phi(y)|0\rangle\).
The scalar inverse with unit source is \(G_F=iD_F\).
These definitions fix factors of \(i\) that differ among references.

For spinors use the gamma matrices of Lectures 13--16,
\(\bar\Psi=\Psi^\dagger\gamma^0\),
\(S^{\mu\nu}=\frac i4[\gamma^\mu,\gamma^\nu]\), and
\(\gamma_*=-i\gamma^0\gamma^1\gamma^2\gamma^3\).
Normal ordering is relative to the free particle--antiparticle vacuum.
\end{conventionbox}

The six original prompts are numbered consecutively as Problems 1--6.
Problems 7 and 8 are new. Work in three spatial dimensions except where
Problem 3 explicitly asks for general \(d\).
\begingroup\small\tableofcontents\endgroup

\problem{1}{Retarded and advanced Green functions}
\lectureconnection{17, Sections 1, 4, and 5}
Start from the scalar commutator
\begin{equation}
 \mathcal K(t,\bm r)=\int\dd\Pi_p
 [e^{-iE_{\bm p}t+i\bm p\cdot\bm r}-e^{iE_{\bm p}t-i\bm p\cdot\bm r}].
 \label{eq:K}
\end{equation}
Do not replace the commutator by a retarded function before imposing a
boundary condition.

\begin{parts}
\item Derive \(L\mathcal K=0\), \(\mathcal K(0,\bm r)=0\), and
\(\partial_t\mathcal K(0,\bm r)=-i\delta^{(3)}(\bm r)\).
Check the last sign directly against the canonical equal-time commutator.
\item Differentiate \(G_R=i\theta(t)\mathcal K\) and
\(G_A=-i\theta(-t)\mathcal K\) as distributions and show that both
satisfy \(LG=\delta^{(4)}\). Explain why using the same prefactor for the
two step functions would give opposite sources. Handle the
\(\delta'\mathcal K\) term explicitly or use the jump of the first derivative.
\item Fourier transform only the spatial coordinates and solve the
resulting oscillator equation with the two respective support
conditions. Recover
\[
 G_R(t,\bm p)=\theta(t)\frac{\sin(E_{\bm p}t)}{E_{\bm p}},\qquad
 G_A(t,\bm p)=-\theta(-t)\frac{\sin(E_{\bm p}t)}{E_{\bm p}}.
\]
\item Derive the four-momentum pole prescriptions. Draw the poles in
the complex \(p^0\) plane and show by residues why one answer vanishes
for \(t<0\) and the other for \(t>0\).
\item Couple a smooth source \(J\) of compact spacetime support through
\(J\phi\). Express the solution with no incoming disturbance using
\(G_R\), and prove that it cannot change outside the causal future of
the source. What different boundary condition does \(G_A\) impose?
\item Show that \(G_R-G_A=i\mathcal K\) solves the homogeneous equation.
Explain why the differential equation for a Green function is insufficient
without a boundary condition.
\end{parts}

\problem{2}{The Feynman function as an inverse wave operator}
\lectureconnection{17, Sections 5--6}
For a real scalar field,
\begin{equation}
 D_F(z)=\theta(t)D^+(z)+\theta(-t)D^+(-z),\qquad
 D^+(z)=\int\dd\Pi_p e^{-ipz}.
 \label{eq:DF}
\end{equation}
Keep the distinction between a time-ordered vacuum expectation and a
Green function normalized to a unit delta-function source.

\begin{parts}
\item Apply \(L\) to \eqref{eq:DF} distributionally. Derive
\(LD_F=-i\delta^{(4)}\), using the equal-time jump instead of appealing
to a remembered propagator formula. Hence determine the source equation
for \(G_F=iD_F\).
\item Derive
\[
 D_F(z)=\lim_{\epsilon\downarrow0}\int\frac{\dd^4p}{(2\pi)^4}
 \frac{i e^{-ipz}}{p^2-m^2+i\epsilon}.
\]
Locate both poles and evaluate the \(p^0\) contour separately for
\(t>0\) and \(t<0\), keeping the contour orientation.
\item Compare this pole prescription to those of Problem 1. Compute
\(L(G_F-G_R)\). Is \(G_F\) a causal response kernel? Explain how a
nonzero spacelike \(D_F\) is compatible with microcausality.
\item For a Dirac field, start with fermionic time ordering and the
spin sums to show
\[
 S_F(z)=\langle0|T\Psi(x)\bar\Psi(y)|0\rangle
       =(i\slashed\partial+m)D_F(z).
\]
Explain the minus sign when the time ordering exchanges the two odd fields.
\item Derive the momentum-space numerator and verify
\((i\slashed\partial-m)S_F=i\delta^{(4)}I_4\).
Why would omitting the fermionic time-ordering sign fail to give this
local source equation?
\end{parts}

\problem{3}{Massive two-point functions in arbitrary dimension}
\lectureconnection{15, Sections 2--3; 17, Section 2}
In \(d+1\) dimensions define
\begin{equation}
 D_d^+(t,\bm r)=\lim_{\epsilon\downarrow0}\int
 \frac{\dd^dp}{(2\pi)^d2E_{\bm p}}
 e^{-iE_{\bm p}(t-i\epsilon)+i\bm p\cdot\bm r}.
 \label{eq:Dd}
\end{equation}
Assume integer \(d\geq1\) and \(m>0\). You may use
\begin{align*}
 (p^2+m^2)^{-1/2}&=\frac1{\sqrt\pi}\int_0^\infty
 \dd\alpha\,\alpha^{-1/2}e^{-\alpha(p^2+m^2)},\\
 \int\frac{\dd^dp}{(2\pi)^d}e^{-\alpha p^2+i\bm p\cdot\bm r}
 &=(4\pi\alpha)^{-d/2}e^{-r^2/(4\alpha)},\\
 \int_0^\infty\dd\alpha\,\alpha^{-\nu-1}e^{-a\alpha-b/\alpha}
 &=2(a/b)^{\nu/2}K_\nu(2\sqrt{ab}),\qquad a,b>0.
\end{align*}

\begin{parts}
\item At spacelike separation use a frame with \(t=0\). Evaluate
\eqref{eq:Dd} in terms of \(K_\nu(mr)\), giving both its coefficient and
order \(\nu\) as functions of \(d\). Check the dimensions of the result
using \([\phi]=(d-1)/2\).
\item Verify the explicit answers for \(d=1,2,3\). You may use
\(K_{1/2}(z)=\sqrt{\pi/(2z)}e^{-z}\).
State why these formulas apply at \(r>0\) rather than as ordinary
functions at coincident points.
\item Use \(K_\nu(z)\sim\sqrt{\pi/(2z)}e^{-z}\) to derive the leading
large-distance behavior for general \(d\). Identify the exponential
length scale and the power-law prefactor.
\item Continue the spacelike formula with
\(r\mapsto\rho_\epsilon=\sqrt{r^2-(t-i\epsilon)^2}\), on the branch
continued from real positive \(r\). Explain why the sign of the
\(i\epsilon\) matters at timelike separation. Independently derive the
large positive-time behavior at \(\bm r=0\) by a Gaussian stationary-phase
approximation about \(\bm p=0\), keeping its phase and power of \(t\).
\item For \(d>1\), find the massless limit at nonzero spacelike distance.
Why does \(d=1\), where \(K_0(z)\sim-\log(z/2)-\gamma_E\), require a
separate infrared prescription for the undifferentiated massless field?
\end{parts}

\problem{4}{The commutator, its light-cone terms, and locality}
\lectureconnection{17, Sections 1--4}
Use the boundary values from Problem 3 to calculate
\(\mathcal K_d(z)=D_d^+(z)-D_d^+(-z)\).
This problem concerns a distribution, so a formula that agrees away from
the light cone is not yet a complete answer.

\begin{parts}
\item Show directly from the invariant mass-shell integral that
\(\mathcal K_d\) is odd under \(z\mapsto-z\), solves the homogeneous
Klein--Gordon equation, and vanishes at spacelike separation.
Write its timelike expression as a difference of the two boundary
values of the Bessel formula, retaining the time-orientation information.
\item In \(3+1\) dimensions, evaluate the discontinuity, including the
light-cone distribution, and determine the constants \(A,B\) in
\[
 \mathcal K(t,r)=-i\operatorname{sgn}(t)\left[
 A\,\delta(t^2-r^2)+B\,\theta(t^2-r^2)
 \frac{mJ_1(m\sqrt{t^2-r^2})}{\sqrt{t^2-r^2}}\right].
\]
Check the constants against the initial derivative
\(-i\delta^{(3)}(\bm r)\). You may use
\((x\pm i0)^{-1}=\operatorname{PV}(1/x)\mp i\pi\delta(x)\).
\item Take \(m\to0\) in three spatial dimensions and determine whether
the support is on or inside the light cone. Compare with the massless
retarded wave equation in one spatial dimension by solving its
initial-value problem. Is propagation only on the cone universal in
all dimensions?
\item For a complex scalar, calculate
\([\Phi(x),\Phi^\dagger(y)]\) and \([\Phi(x),\Phi(y)]\) from the two mode
families. Identify where the particle--antiparticle cancellation occurs.
What remains if the creation term is discarded?
\item Let \(A\) and \(B\) be local Hermitian observables in spacelike
separated regions. Prove that applying \(e^{-i\lambda B}\) cannot change
\(\langle A\rangle\), even when \(\langle AB\rangle\) does not factorize.
Explain the distinction between vacuum correlation and signaling.
\end{parts}

\problem{5}{Charge conjugation of a quantized Dirac field}
\lectureconnection{13, Section 6; 16, Sections 3 and 5}
Use
\[
 C=\diag(E,-E)=i\gamma^0\gamma^2,\qquad
 E=i\sigma_2,\qquad \Psi^c=C\bar\Psi^T.
\]
Distinguish this numerical matrix from the quantum symmetry operator
\(\Dc\). The field has the mode expansion
\begin{equation}
 \Psi(x)=\sum_s\int\dd\Pi_p
 [b_s(\bm p)u_s(p)e^{-ipx}+d_s^\dagger(\bm p)v_s(p)e^{ipx}],
 \label{eq:Diracmodes}
\end{equation}
with the covariant CAR of Lecture 16. Choose an invariant vacuum and
spinor phases satisfying \(v_s=C\bar u_s^T\).

\begin{parts}
\item Verify \(C^T=-C\), \(C^{-1}=-C\), and
\(C(\gamma^\mu)^TC^{-1}=-\gamma^\mu\). Show that the charge-conjugate of
a free solution satisfies the same real-mass Dirac equation.
Why is the minus sign in the gamma-matrix identity essential?
\item Show \((\Psi^c)^c=\Psi\). For the explicit wavefunctions of
Lecture 14, determine the relation between their two-component bases
\(\xi_s\) and \(\eta_s\) required by \(v_s=C\bar u_s^T\).
\item Impose \(\Dc\Psi\Dc^{-1}=\Psi^c\). Derive the transformations of
\(b_s,d_s\), their adjoints, and normalized one-particle states.
Indicate how the answer changes under a rephasing or a change of spin basis.
\item Transform the normal-ordered Hamiltonian and charge. Then
establish the transformation of \(:\bar\Psi\Psi:\) and
\(:\bar\Psi\gamma^\mu\Psi:\), keeping the fermionic reordering signs.
\item Explain why \(\Psi\mapsto C\bar\Psi^T\) is antilinear on a
classical complex spinor but \(\Dc\) is unitary on Fock space.
Contrast with time reversal, including its action on \(i\), energy,
momentum, and spin.
\item Add a minimal coupling \(D_\mu=\partial_\mu+iqA_\mu\).
Determine the accompanying transformation of \(A_\mu\) required for a
charge-conjugation symmetry of the full system. Why is transforming
only the spinor in a fixed generic background not a symmetry statement?
\end{parts}

\problem{6}{Gordon identities with fixed spin conventions}
\lectureconnection{14, Sections 6--7; 16, Section 3}
Let \(p^2=p'^2=m^2\), \(p^0,p'^0>0\), and \(m>0\). Define
\(q=p'-p\), \(P=p'+p\), and
\(S^{\mu\nu}=\frac i4[\gamma^\mu,\gamma^\nu]\).
The spin labels \(s,s'\) need not agree. Use
\[
 (\slashed p-m)u_s(p)=0,\qquad
 (\slashed p+m)v_s(p)=0,
\]
and their adjoint equations. If another source uses
\(\sigma^{\mu\nu}=\frac i2[\gamma^\mu,\gamma^\nu]\), then
\(\sigma^{\mu\nu}=2S^{\mu\nu}\).

\begin{parts}
\item Derive \(\gamma^\mu\gamma^\nu=\eta^{\mu\nu}-2iS^{\mu\nu}\).
Using the external equations of motion, express
\(2m\bar u_{s'}(p')\gamma^\mu u_s(p)\) in terms of \(P^\mu\),
\(S^{\mu\nu}\), and \(q_\nu\). Derive the corresponding vector identity
for two \(v\) spinors and explain its mass-sign difference.
\item Starting from
\(\bar u(p')[\slashed p'\gamma^\mu\gamma_*-
\gamma^\mu\gamma_*\slashed p]u(p)=0\), prove
\begin{equation}
 \bar u_{s'}(p')\left[P^\mu+2iS^{\mu\nu}q_\nu\right]
                        \gamma_*u_s(p)=0.
 \label{eq:axialzero}
\end{equation}
Prove the same vanishing identity with both \(u\)'s replaced by \(v\)'s.
Keep the ordering of \(\gamma_*\) explicit.
\item Use the sum, rather than the difference, of the two external
Dirac equations to derive the axial Gordon identity for
\(2m\bar u(p')\gamma^\mu\gamma_*u(p)\).
Express its right side using \(q^\mu\) and \(S^{\mu\nu}P_\nu\).
What changes for the \(v\)-spinor identity?
\item Contract the vector and axial currents with \(q_\mu\).
Derive their on-shell divergences. Explain the different role of the
mass in vector-current and axial-current conservation for the free theory.
\item Check the formulas at \(p'=p\), and then at a nonzero momentum
transfer using the explicit square-root spinors from Lecture 14.
Would a sign error in front of \(S^{\mu\nu}q_\nu\) be detected by the
zero-transfer check alone? Explain how the identities change if the
momentum-transfer convention is reversed.
\end{parts}

\problem{7}{Scalar normalization, vacuum energy, and charge}
\lectureconnection{15, Sections 2--6}
\textit{New problem.}
Quantize a complex scalar with
\(\Lag=\partial_\mu\Phi^\dagger\partial^\mu\Phi-m^2\Phi^\dagger\Phi\).
Use a periodic box of volume \(V=L^3\) when evaluating coincident
momentum delta functions, and a spatial cutoff \(|\bm p|<\Lambda\) when
summing vacuum energies.

\begin{parts}
\item Derive the two canonical momenta and write all nonzero
equal-time commutators. Begin with the covariant expansion
\[
 \Phi(x)=\int\dd\Pi_p[a(\bm p)e^{-ipx}+b^\dagger(\bm p)e^{ipx}].
\]
Determine both mode commutators and verify the canonical brackets.
Explain why the \(1/(2E)\) expansion and a mode commutator without
\(2E\) cannot be combined.
\item Define normalized box oscillators and construct a state with
\(n\) particles in one mode and \(\bar n\) antiparticles in another.
Calculate its energy, charge, and total particle number.
Then describe a normalized one-particle wave packet in the continuum.
\item Derive the un-normal-ordered Hamiltonian and Noether charge.
After normal ordering, show that the Hamiltonian contains the sum of
the two number densities and the charge their difference. Verify
\([Q,\Phi]=-\Phi\) in the phase convention \(\delta\Phi=-i\alpha\Phi\).
\item Calculate the leading \(\Lambda^4\) divergence of the vacuum
energy density for one real scalar and for one complex scalar. Separate
the volume dependence from the ultraviolet dependence. What is the
relation between \((2\pi)^3\delta^{(3)}(0)\) and the box volume?
\item Add \(-\lambda(\Phi^\dagger\Phi)^2\) to the Lagrangian. Which
continuous charge remains conserved? Why need total particle number
not remain conserved? State the energy-momentum constraint on pair
production and why charge conservation alone does not make the vacuum
unstable.
\item Explain the limited role of normal ordering in this example.
What additional information is needed to compare predictions obtained
with two different ultraviolet regulators in an interacting theory?
\end{parts}

\problem{8}{Fermionic quantization, positivity, and local observables}
\lectureconnection{16, Sections 1--4; 17, Section 3}
\textit{New problem.}
Use the free Dirac mode expansion \eqref{eq:Diracmodes}, with independent
\(b,d\) CAR and vacuum annihilated by both families. You may use
\(\sum_su_s\bar u_s=\slashed p+m\) and
\(\sum_sv_s\bar v_s=\slashed p-m\).

\begin{parts}
\item For one normalized mode \(c\), show that \(n=c^\dagger c\) obeys
\(n^2=n\). For two distinct modes, exhibit the exchange sign of their
two-particle state. Explain why \((c^\dagger)^2\vac=0\) is an exclusion
statement, not the statement that all two-fermion states vanish.
\item Derive \(\{\Psi_\alpha(t,\bm x),\Psi_\beta^\dagger(t,\bm y)\}
=\delta_{\alpha\beta}\delta^{(3)}(\bm x-\bm y)\).
Display the change \(\bm p\mapsto-\bm p\) in the antiparticle integral
and the cancellation of the mass and spatial gamma terms.
\item Starting with \(H=\int\dd^3x\,\Psi^\dagger h_D\Psi\), derive
\[
 H=\sum_s\int\dd\Pi_p E_{\bm p}(b_s^\dagger b_s-d_sd_s^\dagger).
\]
Normal-order it and \(Q=\int\dd^3x:\Psi^\dagger\Psi:\).
Verify directly the energy and charge of \(b_s^\dagger\vac\) and
\(d_s^\dagger\vac\).
\item Temporarily replace both CAR families by positive-norm bosonic
oscillators in the preceding Hamiltonian. Show why the energy cannot
be bounded below. Explain why a negative mode commutator is not a
positive-Hilbert-space remedy with the same annihilation vacuum.
\item Derive
\[
 \{\Psi_\alpha(x),\bar\Psi_\beta(y)\}
 =(i\slashed\partial+m)_{\alpha\beta}\mathcal K(x-y)
\]
and establish its spacelike vanishing. For separated points, show why
two even bilinears such as \(:\bar\Psi\Psi:\) commute outside the light
cone. State the role of a composite-operator prescription at coincidence.
\item At a finite regulator, set \(f_s=d_s^\dagger\) and reinterpret the
vacuum as a filled set of negative-energy modes. Derive the energy and
charge of a hole relative to that filled state. Which parts of this
argument require care when the number of modes becomes infinite?
\end{parts}
\end{document}
